\documentclass[../../main.tex]{subfiles}
\begin{document}

\begin{flushright}
\footnotesize\textit{【自動文字起こし・要確認】}
\end{flushright}

\noindent
{\bf [解]} $X=m, Y=n$ である事象を $X_m, Y_n$ とおく．$X_m, Y_n$ が独立だから
\begin{align}
P(X_m \wedge Y_n) = P(X_m) P(Y_n) \quad \dots \text{①}
\end{align}
又，$X+Y=n$ となる確率が $(n+1) P_{n+1}$ に等しく，この時 $(X, Y) = (k, n-k) \ (k=0, 1, \dots, n)$ だから
\begin{align}
& P(X+Y=n) = (n+1) P_{n+1} \\
\therefore & \sum_{k=0}^n P(X_k \wedge Y_{n-k}) = (n+1) P_{n+1} \quad \dots \text{②}
\end{align}
①を②に代入して
\begin{align}
\sum_{k=0}^n P_k P_{n-k} = (n+1) P_{n+1} \quad \dots \text{③}
\end{align}
以下，任意の $n \in \mathbb{Z}_{\ge 0}$ に対し，$P_n = P_0^{n+1} \dots \text{◆}$ であること $\dots$ ④ を帰納的に示す． \\
$n=0$ の時は明らかに成立する．以下 $n \le l \ (l \in \mathbb{Z}_{\ge 0})$ での◆の成立を仮定する． \\
③から
\begin{align}
(l+1) P_{l+1} = \sum_{k=0}^l P_0^{k+1} P_0^{l-k+1} = (l+1) P_0^{l+2}
\end{align}
$l+1 \neq 0$ から
\begin{align}
P_{l+1} = P_0^{l+2}
\end{align}
となり，$n=l+1$ でも◆は成立．以上から④は示された．これを $\sum_{n=0}^\infty P_n = 1$ に代入する． \\
$P_0 = 1$ は明らかに不適なので，$P_0 \neq 1$ とすると，
\begin{align}
\sum_{n=0}^N P_n = P_0 \dfrac{1 - P_0^{N+1}}{1 - P_0} \quad \dots \text{⑤}
\end{align}
これが収束するのは，$|P_0| < 1$ の時で，このもとで⑤は $\dfrac{P_0}{1-P_0}$ に収束する．
\begin{align}
\dfrac{P_0}{1-P_0} = 1 \iff P_0 = \dfrac{1}{2} \quad (|P_0| < 1 \text{をみたす})
\end{align}
だから◆に代入して
\begin{align}
P_n = \Bigl(\dfrac{1}{2}\Bigr)^{n+1}
\end{align}
だから，$r = \dfrac{1}{2}$ として
\begin{align}
\sum_{k=0}^n k \cdot P_k &= \dfrac{1}{1-r} r - \dfrac{1}{(1-r)^2} r^{n+1} + \dfrac{r}{1-r} \\
&\xrightarrow{n \to \infty} \dfrac{r}{1-r} = 1 \quad (|r| < 1)
\end{align}
となる．

\newpage

\end{document}
